
\def\PUDcodigo{ECA230}

\def\PUDcodacad{04507.57}

\def\PUDdisciplina{Sistemas Embarcados}

\def\PUDsemestre{Opt}

\def\PUDhoras{80}

%NOVOS ---------------------------------------------------
\def\PUDhorasTeoria{40}
\def\PUDhorasPratica{40}

\def\PUDinterECA{ 
\hyperref[PUD:ECA213]{Microcontroladores (04507.28)}

\hyperref[PUD:ECA202]{Controle II (04507.40)}

\hyperref[PUD:ECA225]{Dispositivos Periféricos (04507.42)} 
}

\def\PUDatuaisECA{
- Tendências modernas de sistemas embarcados, como Internet das Coisas (IoT) 

- Aplicações de sistemas embarcados na indústria moderna (4.0) e no dia-a-dia.
}

\def\PUDrecursos{Projetor multimídia; Quadro branco e pincel.}

%----------------------------------------------------------

\def\PUDcreditos{4}

\def\PUDprerequisito{ \hyperref[PUD:ECA225]{ECA225} }

\def\PUDpreacad{ \hyperref[PUD:ECA225]{Dispositivos Periféricos (04507.42)} }

\def\PUDobjetivos{Compreender, projetar e desenvolver sistemas embarcados e sistemas em tempo real.}

\def\PUDementa{Introdução a sistemas em tempo real; 
Confiabilidade e tolerância a falhas;  
Programação concorrente; 
Comunicação e sincronização baseada em memória compartilhada;  
Sincronização baseada em mensagem;  
Ações atômicas e processos concorrentes;  
Controle de recurso;  
Facilidades em tempo real; 
Escalonamento adaptativo;  
Protocolos de comunicação, sistemas operacionais e middleware de tempo real; 
Entrada e saída de dados em sistemas embarcados;  
Desenvolvimento de sistemas de tempo real;  
Sistemas operacionais para sistemas embarcados;  
Ferramentas de desenvolvimentos para sistemas embarcados; 
Linguagens de programação para sistemas embarcados; 
Plataformas de hardware para sistemas embarcados;  
Projeto e desenvolvimento de sistemas embarcados.}

\def\PUDconteudo{ & Sistemas em tempo real: Definições, características e exemplos de sistemas em tempo real. Confiabilidade e tolerância a falhas: confiabilidade, falha, falta e erro. Prevenção de falhas e tolerância a falhas;  programação N-versões.  Redundância dinâmica de software. Bloco de recuperação para tolerância a faltas de software. Programação concorrente: noções de processo;  execução concorrente.  Representação de processos. Sistema em tempo-real  simples. Comunicação e sincronização baseada em memória compartilhada: exclusão mútua e condição de sincronização. Sincronização e comunicação baseada em mensagem. Ações atômicas e processos concorrentes: ações atômicas e estas em linguagem concorrentes. Controle de recurso: controle de recursos e ações atômicas, gerenciamento de recursos, potência expressiva e facilidade de uso, uso de recurso, deadlock.  
\\ 
 & Sistemas embarcados: Sistemas operacionais para sistemas embarcados: Windows CE, microlinux, Android e IOS. Ferramentas de desenvolvimento: eclipse, linguagem c/c++, java, xml. Plataformas de hardware: processadores, ARM, microcontroladores. Plataformas de desenvolvimento: pic, dsPic, Arm, MSP e smartphones.  
\\ 
 & Aplicações práticas de sistemas embarcados e sistemas em tempo real. Projeto e desenvolvimento de sistemas embarcados. }

\def\PUDmetodo{Aulas expositórias que podem ser teóricas e/ou práticas, onde as práticas no laboratório serão marcadas no decorrer da disciplina.}

\def\PUDavalia{A avaliação será feita com aplicação de provas teóricas e/ou práticas, além da possibilidade de inclusão de trabalhos e seminários no decorrer da disciplina.}

\def\PUDbibbasica{ &  \href{www.biblioteca.ifce.edu.br/index.asp?codigo_sophia=24005}{Oliveira}, André Schneider de.  Sistemas embarcados: hardware e firmware na prática.  São Paulo: Érica.  2006.  ISBN: 8536501057.
\\ 
 &  \href{www.biblioteca.ifce.edu.br/index.asp?codigo_sophia=62221}{Shaw}, Alan C.  Sistemas e software de tempo real.  Porto Alegre: Bookman, 2003.  240p.  ISBN: 8536301724.
\\ 
 & \href{www.biblioteca.ifce.edu.br/index.asp?codigo_sophia=24095}{Yaghmour}, K;  Masters, J.;  Ben-Yossef, G.;  Gerum, P.  Construindo sistemas linux embarcados.  2º ed.  Rio de Janeiro: Alta Books.  2009.  377p.  ISBN: 9788576083436. }

\def\PUDbibcomplementar{ &  \href{www.biblioteca.ifce.edu.br/index.asp?codigo_sophia=24064}{Tanenbaum}, Andrew S.  Sistemas operacionais modernos.  2º ed.  São Paulo, SP : Pearson Prentice Hall, 2007.  695p.  ISBN: 9788587918574.
\\
 &  \href{www.biblioteca.ifce.edu.br/index.asp?codigo_sophia=39627}{Sommerville}, Ian, A.  Engenharia de Software.  9º ed.  E-book, 2011. 529p. ISBN: 9788579361081.
 \\
 &  \href{www.biblioteca.ifce.edu.br/index.asp?codigo_sophia=23899}{Sousa}, Daniel Rodrigues de.  Microcontroladores ARM7 (Philips - Família LPC213X): o poder dos 32 bits: teoria e prática.  São Paulo, SP: Érica, 2006.  278p.  ISBN: 8536501200.
\\
 &  \href{www.biblioteca.ifce.edu.br/index.asp?codigo_sophia=46123}{Bräunl}, Thomas.  Embedded robotics: mobile robot design and applications with embedded systems.  3º ed.  Perth, Austrália: Springer, 2008.  541p.  ISBN: 9783540705338.
\\
 &  \href{www.biblioteca.ifce.edu.br/index.asp?codigo_sophia=57239}{Tanenbaum}, Andrew S.  Organização estruturada de computadores. E-book. 6º ed.  Pearson.  628p.  ISBN: 9788581435398. }
%  BURNS, A.  WELLINGS, A.  Real-Time Systems and Programming Languages.  2nd edition.  Addison: Wesley, 1997.  ISBN 9780201175295. 
%\\ 
 %&  HERMAN,  Kopetz;  “Real-Time  Systems: design principles for distributed embedded applications;  Kluwer Academic Publishers, 1997.  ISBN 9780792398943
%\\ 
 %&  LI, Jane.  Real-Time Systems.  Prentice-Hall, 2000.  ISBN 9780130996510
%\\ 
 %&  BUTTAZZO, G.  Hard Real-Time Computing Systems - Predictable Scheduling Algorithms and Applications.  Kluwer Academic Publishers, 199
%\\ 
 %&  FARINES, J. M.  FRAGA, J.  da S.  OLIVEIRA, R.  S.  Sistemas de Tempo Real.  Escola de Computação 2000, IME-USP, São Paulo-SP, julho/2000.

\def\PUDautor{Pedro Pedrosa Rebouças Filho}

\def\PUDdata{2013-06-06}

\def\PUDrevisao{2}

\def\PUDrevisor{Luiz Daniel Santos Bezerra}

\def\PUDdatarevisao{2019-05-15}

\PUDcorpo
